\chapter*{Wprowadzenie i zakres}
\phantomsection
\addcontentsline{toc}{chapter}{Wprowadzenie i zakres}
Niniejszy skrypt wyodrębnia z wcześniejszej książki część poświęconą
metodom topologii algebraicznej. Zachowano zasadniczą treść rozdziałów
źródłowych: twierdzenia, dowody, przykłady i rysunki.
Dodano osobne wprowadzenie topologiczne, samodzielny rozdział
z algebry abstrakcyjnej, kohomologię singularną,
fibracje, teorię \v{C}echa, ciągi spektralne, przestrzenie klasyfikujące, teorię przeszkód, wieże Postnikowa, operacje kohomologiczne, homotopię stabilną, konstrukcję Pontriagina--Thoma, homologię Morse'a, twierdzenie Lefschetza, Poincarégo--Hopfa, homologię Borel--Moore'a, kohomologię snopów, dualność Verdiera, formalizm sześciu operacji i cykle pobliskie oraz topologię odmian projektowych, twierdzenia Lefschetza i rozdmuchania. Rozdział 29 uzupełnia granice odwrotne, teorię przeszkód,
ekscyzję homotopijną i analizę Morse'a.
Rozdział 30 stanowi zwięzły dodatek z geometrii różniczkowej.
Część poświęcona dualności i klasom charakterystycznym pochodziła z
późniejszej części skryptu geometrycznego.

Punktem wyjścia są podstawy topologii i algebry liniowej;
potrzebną teorię grup, pierścieni i modułów przedstawiamy od początku.
W rozdziale o formach i w rozdziale o klasach charakterystycznych
korzystamy również z elementarnej teorii gładkich rozmaitości,
przestrzeni stycznych, wiązek wektorowych i transwersalności;
przypomnienie potrzebnych pojęć znajduje się w dodatku.
Pełne dowody twierdzeń o otoczeniach tubularnych, transwersalności
oraz teorii spinowej pozostają w książce o geometrii różniczkowej.

\noindent\textbf{Status dowodów.} Słowo „dowód” oznacza miejscami
dowód \emph{względny}, to znaczy przeprowadzony po przyjęciu
wyraźnie wskazanego twierdzenia zewnętrznego,
na przykład o transwersalności, rozkładzie lub dualności Verdiera.
Uzasadnienia takie nie stanowią niezależnych dowodów
przyjmowanych twierdzeń. Ich zakres i zależności
należy odczytywać zgodnie z uwagami zawartymi w rozdziałach.
Nie utożsamiamy szkiców i rachunków wynikających z głębokiego
twierdzenia z dowodem tego twierdzenia od podstaw.

\noindent\textbf{Uwaga redakcyjna.} Ta wersja scala istniejący materiał;
nie oznacza pełnej przebudowy dowodów wszystkich dużych twierdzeń.
Przykłady i konstrukcje zachowano, a ich ewentualne powtórzenia
w innych tomach nie są automatycznie usuwane bez porównania
założeń, wersji wyników i istotnych szczegółów dowodowych.
Osobny dokument z zadaniami zawiera wskazówki do trudniejszych ćwiczeń.
\clearpage
